\chapter{Considerações Finais}
\label{cap:consideracoes}

Este trabalho partiu de uma constatação simples e, ao mesmo tempo, incômoda: os sistemas do
Projeto EMA, embora tenham cumprido os seus objetivos de pesquisa, não podem hoje ser
reproduzidos por um novo integrante sem esforço considerável. A pergunta que orientou a
investigação --- de que forma o instrumental da Engenharia de Software pode ser aplicado ao
diagnóstico, à documentação e à reestruturação desses sistemas --- foi respondida por meio de
dois estudos de caso complementares, o sistema de realidade virtual e o sistema de ciclismo
com eletroestimulação funcional, examinados sob uma postura de consultoria.

\section{Contribuições}
\label{sec:contribuicoes}

As principais contribuições deste trabalho são de natureza diagnóstica e metodológica.
Documentou-se e diagramou-se a arquitetura atual de dois sistemas até então pouco descritos,
reconstruindo-a por análise de código e engenharia reversa. Mapeou-se, de forma sistemática,
a lacuna entre o que está \emph{registrado} e o que está de fato \emph{montado} --- desde a
documentação perdida na migração de ROS~1 para ROS~2, no caso do ciclismo, até a lógica
embarcada em \emph{Blueprints} e a coexistência de mecanismos de comunicação, no caso da
realidade virtual. A partir desse diagnóstico, consolidaram-se recomendações de Engenharia de
Software organizadas por prioridade --- versionamento, documentação, diagramação,
modularização e caminho de integração ---, que constituem um roteiro acionável para as etapas
seguintes. O conjunto entrega ao laboratório um retrato fiel do seu ponto de partida e um
plano para evoluir a partir dele.

\section{Limitações do Trabalho}
\label{sec:limitacoes-trabalho}

Reconhecem-se as limitações desta etapa. Consoante o seu caráter preparatório, o escopo foi o
de diagnóstico, documentação e proposta --- não o da implementação e da validação
experimental das soluções, reservadas à etapa subsequente. Ademais, o prazo de execução não
permitiu a reprodução integral de sistemas cuja documentação é ausente ou incompleta; essa
própria impossibilidade, contudo, não é uma omissão, mas um dos resultados diagnósticos do
trabalho, e reforça a pertinência das recomendações apresentadas.

\section{Trabalhos Futuros}
\label{sec:trabalhos-futuros}

Como continuidade natural, aponta-se a implementação das recomendações do
Capítulo~\ref{cap:proposta}: a padronização dos repositórios e da documentação, a
reativação efetiva do sistema de realidade virtual, a reestruturação e a modularização do
sistema de ciclismo com FES e, ao final, a convergência das duas frentes em um ecossistema de
reabilitação integrado e \emph{plug-and-play}. Consolidada essa base de software ---
padronizada, documentada e resiliente ---, o laboratório estará apto a retomar, com
segurança e continuidade, a evolução de tecnologias cujo propósito último é ampliar a
autonomia e a qualidade de vida de pessoas em reabilitação.
